\setcounter{subsection}{0}
\section{CONSTRUCCIÓN E IMPLEMENTACIÓN DEL ARTEFACTO}
\sloppy

La construcción del artefacto materializó la Fase 4 (\textit{Desarrollo}) de la metodología de Investigación en Ciencia del Diseño (\textit{Design Science Research}, DSR) formulada por Peffers et al. (2007). A partir de la arquitectura hexagonal (Cockburn, 2024) y los contratos unificados definidos durante la Fase 3 en el Documento de Arquitectura de Software (SAD), la implementación se ejecutó sobre un monorrepositorio en TypeScript bajo modo estricto (\texttt{strict: true} en \texttt{tsconfig.json}) y Node.js 20 LTS. De acuerdo con el Plan de Gestión del Proyecto de Software (SPMP), el artefacto construido se compone de tres elementos principales integrados en el mismo árbol de código: el SDK unificado (\texttt{sdk/}), la API de Simulación con su fachada REST (\texttt{simulator-api/}) y el catálogo verificable de datos de prueba por pasarela (\texttt{docs/testing-data/}), complementados con los ejemplos ejecutables (\texttt{examples/}) y los cuatro prototipos de comercio electrónico (\texttt{prototypes/}). Las decisiones técnicas que modificaron contratos, umbrales o fronteras entre módulos quedaron registradas de forma cronológica en la bitácora de arquitectura del proyecto, cuyos puntos se referencian a lo largo de este capítulo.

\subsection{Proceso de desarrollo}

\subsubsection{Iteraciones de construcción y trazabilidad}

El desarrollo se dividió en tres iteraciones secuenciales de dos semanas cada una, ejecutadas entre el 25 de agosto y el 31 de octubre de 2026, siguiendo la planificación establecida en el SPMP. En cada iteración, los requisitos funcionales y no funcionales se desglosaron en tareas trazables vinculadas a ramas temáticas e integradas mediante solicitudes de cambio (\textit{pull requests}) revisadas por pares.

La \textbf{Iteración 1} (\textit{Núcleo del SDK}) estableció el esqueleto del paquete \texttt{sdk/}, la entidad inmutable \texttt{Transaction} con preservación del estado nativo en \texttt{rawStatus} (punto 3 de la bitácora de arquitectura), los objetos de valor de dominio (\texttt{Amount}, \texttt{PaymentMethod}, \texttt{PaymentResult}, \texttt{TaxBreakdown}, \texttt{Credentials} y \texttt{WebhookEvent}), la jerarquía de errores \texttt{KitPagosError} con sus trece códigos tipados, el puerto de salida \texttt{PaymentGatewayPort}, los servicios de aplicación (\texttt{ResponseNormalizer}, \texttt{RetryHandler} y \texttt{ErrorHandler}) y la fachada pública \texttt{KitPagos} junto con \texttt{SdkConfigurator} y \texttt{GatewayFactory}, aplicando los patrones \textit{Facade}, \textit{Factory Method} y \textit{Adapter} (Gamma et al., 1994). Durante esta iteración se fijó también la regla de marcadores explícitos para métodos en construcción (\texttt{throw new Error("aun no esta implementado")}), evitando que retornos vacíos o valores nulos se confundieran con operaciones completadas (punto 6).

La \textbf{Iteración 2} (\textit{Adaptadores de pasarela}) implementó los cuatro adaptadores de infraestructura en \texttt{sdk/src/\allowbreak infrastructure/\allowbreak adapters/} (\texttt{WompiAdapter}, \texttt{MercadoPagoAdapter}, \texttt{KushkiAdapter} y \texttt{RapydAdapter}) conforme a las especificaciones técnicas de cada proveedor (Kushki, 2024; Mercado Pago, 2024; Rapyd Financial Technology, 2024; Wompi Colombia, 2024), los manejadores de verificación criptográfica de webhooks en tiempo constante (\texttt{WebhookVerifier} y los cuatro verificadores por proveedor en \texttt{sdk/src/\allowbreak domain/\allowbreak services/}) y la batería de pruebas de contrato contra los sandboxes reales de las pasarelas (\texttt{sdk/test/\allowbreak sandbox/}, puntos 43 a 58). Los hallazgos empíricos obtenidos al ejecutar cobros reales en los entornos de prueba de los proveedores retroalimentaron el dominio: motivaron la introducción del resultado discriminado \texttt{PaymentResult} con \texttt{PendingRedirect} para cobros que exigen redirección del pagador (punto 39), la aceptación de un mapa de URLs base por pasarela en \texttt{SdkConfigurator} (punto 57) y la separación del estado nativo de tarjeta y de transferencia en Kushki (punto 48).

La \textbf{Iteración 3} (\textit{API de simulación y soporte de validación}) se organizó en cuatro entregables internos conforme al plan de trabajo del proyecto: (3.1) las rutas de imitación de pasarela bajo \texttt{/v1/sim/:gateway/*} y la fachada REST unificada bajo \texttt{/v1/api/*} en \texttt{simulator-api/} siguiendo las restricciones del estilo arquitectónico REST (Fielding, 2000), construidas sobre Fastify con almacenamiento en memoria por recurso y máquinas de estado declarativas (puntos 7 a 13 y 63); (3.2) el emisor manual y automático de webhooks firmados junto con el catálogo determinista de escenarios controlados por encabezados, credenciales, datos de prueba y montos reservados (puntos 12 y 83); (3.3) la documentación técnica por módulo, la especificación interactiva OpenAPI en \texttt{/docs} (punto 125), la colección versionada de Postman (\texttt{simulator-api/\allowbreak postman/}) y el despliegue público en Render (\texttt{render.yaml}, puntos 79, 84 y 88); y (3.4) el módulo de tokenización en navegador \texttt{kit-pagos-colombia/\allowbreak browser} (\texttt{sdk/src-browser/}, puntos 77, 78, 81 y 106), los doce ejemplos ejecutables en \texttt{examples/} y los cuatro prototipos de comercio electrónico en \texttt{prototypes/} (puntos 91 a 124).

En la Tabla \ref{tab:iteraciones-fase4} se sintetizan las tres iteraciones de la Fase 4 junto con sus entregables verificables en el repositorio.

\begin{table}[H]
\centering
\small
\caption{Iteraciones de construcción ejecutadas durante la Fase 4 (Desarrollo)}
\label{tab:iteraciones-fase4}
\begin{tabular}{|p{2.6cm}|p{4.3cm}|p{6.9cm}|}
\hline
\textbf{Iteración} & \textbf{Alcance principal} & \textbf{Artefactos y directorios verificables} \\
\hline
Iteración 1: Núcleo del SDK & Entidad \texttt{Transaction}, objetos de valor, puerto \texttt{PaymentGatewayPort}, normalización, reintentos, manejo de errores y fachada \texttt{KitPagos}. & \texttt{sdk/src/domain/}, \texttt{sdk/src/\allowbreak application/}, \texttt{sdk/src/\allowbreak infrastructure/\allowbreak KitPagos.ts}, \texttt{SdkConfigurator.ts}, \texttt{GatewayFactory.ts}. \\
\hline
Iteración 2: Adaptadores de pasarela & Adaptadores para Wompi, Mercado Pago, Kushki y Rapyd; verificadores HMAC/SHA-256 de webhooks; pruebas de contrato en sandboxes. & \texttt{sdk/src/\allowbreak infrastructure/\allowbreak adapters/}, \texttt{sdk/src/\allowbreak domain/\allowbreak services/}, \texttt{sdk/test/\allowbreak sandbox/}, \texttt{docs/testing-data/}. \\
\hline
Iteración 3: API de simulación y validación & Superficies \texttt{/v1/sim} y \texttt{/v1/api}, escenarios deterministas, módulo \texttt{kit-pagos-colombia/\allowbreak browser}, despliegue en Render, ejemplos y prototipos A, B, C y D. & \texttt{simulator-api/}, \texttt{sdk/src-browser/}, \texttt{examples/}, \texttt{prototypes/}, \texttt{render.yaml}, \texttt{simulator-api/\allowbreak postman/}. \\
\hline
\end{tabular}
\end{table}

\subsubsection{Integración continua y guardas automatizadas}

El flujo de integración continua se codificó en \texttt{.github/\allowbreak workflows/\allowbreak ci.yml} y se ejecuta en cada \textit{push} y en cada \textit{pull request} hacia las ramas de integración (\texttt{main}, \texttt{develop}, \texttt{testing}, \texttt{devops}, \texttt{feature/**} y \texttt{fix/**}). El flujo está dividido en siete trabajos independientes sobre \texttt{ubuntu-latest} con Node.js 20:

\begin{enumerate}
    \item \textbf{\texttt{sdk}}: ejecuta el análisis estático con ESLint (\texttt{npm run lint}), la verificación de formato con Prettier (\texttt{npm run format:check}), la compilación de ambos puntos de entrada y del paquete de navegador (\texttt{npm run build}), la guarda de complejidad y acoplamiento (\texttt{npm run metrics:ck}), la batería de pruebas unitarias e integración con cobertura (\texttt{npm run test:coverage}), la compilación en seco de los fragmentos de código incluidos en la guía principal del paquete (\texttt{npm run check:readme}, puntos 76 y 115) y la verificación del artefacto empaquetado (\texttt{npm run check:published}).
    \item \textbf{\texttt{examples}}: instala las dependencias del SDK y de \texttt{examples/} y verifica que los doce scripts de demostración compilen sin errores de tipos contra la superficie pública exportada por \texttt{kit-pagos-colombia} (\texttt{npx tsc --noEmit}, punto 21).
    \item \textbf{\texttt{simulator-api}}: compila el SDK y la API de Simulación, valida el estilo y formato del código, ejecuta la guarda de métricas CK sobre \texttt{simulator-api/src} con el mapa de excepciones documentadas (\texttt{ck-exceptions.json}, punto 107) y corre la suite de pruebas de la API con cobertura (\texttt{npm run test:coverage}).
    \item \textbf{\texttt{landing}}: verifica el formato y los enlaces de la página estática informativa alojada en \texttt{landing/}.
    \item \textbf{\texttt{prototypes}}: ejecuta en matriz las verificaciones de tipos (\texttt{typecheck}) y las pruebas automatizadas con Vitest de los tres prototipos con servidor propio (\texttt{checkout-con-sdk-solo-back}, \texttt{checkout-con-sdk-front-back} y \texttt{checkout-sin-solucion}, puntos 110, 116 y 123).
    \item \textbf{\texttt{prototype-frontend-rest}}: compila el paquete de navegador del SDK y ejecuta la verificación de tipos, las pruebas y el empaquetado con Vite del prototipo estático \texttt{checkout-frontend-rest} (punto 91).
    \item \textbf{\texttt{prototypes-parity}}: ejecuta el script \texttt{prototypes/\allowbreak check-skeleton-parity.sh}, el cual compara byte a byte todos los archivos de esqueleto compartido entre el prototipo de tratamiento (B) y el prototipo de control (D) fuera de la costura de pagos, garantizando que cualquier divergencia accidental rompa la integración continua antes de fusionarse (punto 116).
\end{enumerate}

Adicionalmente, el repositorio cuenta con dos flujos de trabajo especializados. El archivo \texttt{.github/\allowbreak workflows/\allowbreak publish.yml} automatiza la publicación del paquete en el registro público de npm cuando se crea una etiqueta \texttt{sdk-v*}, exigiendo que la versión de la etiqueta coincida con \texttt{sdk/package.json}, ejecutando todas las guardas de calidad, publicando con procedencia criptográfica (\texttt{--provenance}) en el canal \texttt{next} para versiones preliminares o \texttt{latest} para versiones estables, y verificando la instalación limpia del paquete recién publicado desde el registro remoto (punto 108). Por su parte, \texttt{.github/\allowbreak workflows/\allowbreak final-doc.yml} compila en contenedor TeX Live los documentos oficiales en LaTeX ubicados en \texttt{docs/final-doc/} (punto 126).

La verificación del paquete empaquetado (\texttt{sdk/scripts/\allowbreak check-published-package.ts}) evita que un artefacto incompleto llegue a npm: ejecuta \texttt{npm pack --dry-run --json}, comprueba la presencia de los archivos obligatorios (\texttt{dist/index.js}, \texttt{dist/index.d.ts}, \texttt{dist/browser/index.js}, \texttt{dist/browser/index.d.ts}, \texttt{dist/browser/\allowbreak kit-pagos-browser.iife.js}, la documentación principal del paquete y \texttt{LICENSE}), prohíbe la inclusión de archivos de prueba o mapas de origen (\texttt{*.test.*}, \texttt{*.spec.*}, \texttt{*.map}), instala el paquete en un directorio temporal aislado y ejecuta pruebas de humo que importan el módulo principal y el subcamino \texttt{kit-pagos-colombia/\allowbreak browser} tanto desde CommonJS como desde módulos ES y TypeScript con \texttt{moduleResolution: "bundler"}.

Por su parte, las pruebas de contrato contra los sandboxes reales (\texttt{sdk/test/sandbox/}, ejecutadas mediante \texttt{npm run test:sandbox}) quedaron desacopladas del flujo de CI ordinario porque dependen de la disponibilidad de red y de las credenciales de comercio en los entornos de prueba de Wompi, Mercado Pago, Kushki y Rapyd (punto 43). Siguiendo los principios de diseño de pruebas de software (Ammann \& Offutt, 2016), cuando las variables de entorno de una pasarela están presentes en \texttt{sdk/.env}, el arnés ejecuta las pruebas reales contra el proveedor; cuando faltan, los bloques se marcan como omitidos (\texttt{describe.skipIf}) sin falsificar resultados (punto 46).

\subsubsection{Métricas de Chidamber y Kemerer como guarda de diseño}

Para impedir la degradación estructural del código a medida que se incorporaban particularidades de cada pasarela, se implementó el analizador estático \texttt{sdk/scripts/ck-metrics.ts} utilizando \texttt{ts-morph} sobre el árbol de sintaxis abstracta (AST) de TypeScript, fundamentado en la batería de métricas de diseño orientado a objetos de Chidamber y Kemerer (1994) y en la medida de complejidad ciclomática de McCabe (1976). El analizador evalúa cuatro indicadores sobre cada unidad de diseño en \texttt{sdk/src/} y \texttt{simulator-api/src/}:

\begin{itemize}
    \item \textbf{WMC (\textit{Weighted Methods per Class} $\le 20$)}: suma de la complejidad ciclomática de McCabe (1976) de todos los métodos y constructores de la clase (Chidamber \& Kemerer, 1994). Inicialmente el script contaba únicamente el número de métodos (peso unitario); en el punto 34 de la bitácora de arquitectura se corrigió para sumar las decisiones reales (\texttt{if}, \texttt{else if}, \texttt{for}, \texttt{while}, \texttt{case}, \texttt{catch}, operadores ternarios y cortocircuitos lógicos \texttt{\&\&}, \texttt{||} y \texttt{??}), elevando el umbral de 15 a 20 y añadiendo un tope por método individual.
    \item \textbf{MAX\_CC ($\le 10$)}: complejidad ciclomática máxima permitida en cualquier método o función individual (McCabe, 1976; punto 34).
    \item \textbf{CBO (\textit{Coupling Between Objects} $\le 5$)}: cantidad de tipos propios del proyecto referenciados por la unidad analizada (Chidamber \& Kemerer, 1994). En el punto 107 se refinó el cálculo para recorrer los nodos \texttt{Identifier} del AST cuyos símbolos se resuelven a declaraciones dentro de \texttt{src/}, excluyendo primitivos del lenguaje, tipos globales de Node.js y bibliotecas externas como \texttt{axios} o \texttt{big.js}.
    \item \textbf{RFC (\textit{Response For a Class} $\le 20$)}: suma de los métodos propios de la clase más el conjunto de métodos externos distintos invocados desde sus cuerpos (Chidamber \& Kemerer, 1994).
\end{itemize}

En el punto 87 de la bitácora de arquitectura, el script se extendió con el modo \texttt{--unit file} para evaluar también proyectos basados en módulos y funciones libres de primer nivel —donde cada archivo \texttt{.ts} constituye la unidad de análisis—, criterio que se empleó posteriormente en el protocolo de evaluación de la Fase 5.

Cuando el analizador detectó clases que superaban los umbrales durante la construcción, la respuesta prioritaria fue refactorizar el diseño extrayendo módulos cohesivos: en el punto 33 se extrajeron \texttt{big-arithmetic.ts}, \texttt{minor-units.ts}, \texttt{signature-utils.ts} y \texttt{rapyd-signature.ts}; en el punto 34 se dividió \texttt{ResponseNormalizer} en cuatro normalizadores por pasarela y se separaron los manejadores de webhook y los submódulos de cada adaptador (\texttt{wompi-pse.ts}, \texttt{mercadopago-pse.ts}, \texttt{mercadopago-card.ts}, \texttt{kushki-charge.ts}, \texttt{kushki-pse.ts}, \texttt{kushki-amount.ts}, \texttt{rapyd-checkout.ts}, \texttt{rapyd-pse.ts} y \texttt{rapyd-payload.ts}). Únicamente en los puntos de unión arquitectónica donde reducir el indicador habría obligado a introducir indirecciones artificiales o perder seguridad de tipos se registraron excepciones acotadas en \texttt{KNOWN\_EXCEPTIONS} (\texttt{sdk/scripts/ck-metrics.ts}) y en \texttt{simulator-api/ck-exceptions.json}, resumidas en la Tabla \ref{tab:excepciones-ck}.

\begin{table}[H]
\centering
\small
\caption{Excepciones acotadas de métricas CK registradas en la guarda automatizada}
\label{tab:excepciones-ck}
\begin{tabular}{|p{3.6cm}|p{1.9cm}|p{1.8cm}|p{6.5cm}|}
\hline
\textbf{Clase / Módulo} & \textbf{Métrica y tope} & \textbf{Referencia} & \textbf{Justificación arquitectónica} \\
\hline
\texttt{Transaction} (\texttt{sdk}) & CBO $\le 7$ & Punto 22 & Entidad raíz del dominio que agrega los tipos \texttt{Gateway}, \texttt{TransactionStatus}, \texttt{RejectionReason}, \texttt{Amount}, \texttt{Currency}, \texttt{PaymentMethod} y \texttt{CustomerInfo}. \\
\hline
\texttt{Amount} (\texttt{sdk}) & WMC $\le 23$ & Punto 34 & Objeto de valor financiero con doce métodos e invariantes de cadena decimal canónica, moneda ISO 4217 y desglose tributario \texttt{TaxBreakdown}. \\
\hline
\texttt{KitPagos}, \texttt{GatewayFactory} y los 4 adaptadores (\texttt{sdk}) & CBO $\le 6$ & Punto 47 & Puntos de ensamblaje que implementan \texttt{PaymentGatewayPort} e incorporaron el tipo \texttt{PseBank} al añadir la consulta de bancos PSE (\texttt{getPseBanks}). \\
\hline
\texttt{KitPagosProvider} (\texttt{simulator-api}) & CBO $\le 7$ & Punto 107 & Adaptador entre la superficie HTTP \texttt{/v1/api} y el SDK que coordina credenciales, entorno destino y fábrica de instancias por petición. \\
\hline
\end{tabular}
\end{table}

\subsection{SDK unificado}

\subsubsection{Estructura hexagonal del núcleo}

El paquete \texttt{kit-pagos-colombia} (\texttt{sdk/src/}) consta de 63 archivos de producción distribuidos estrictamente en las tres capas de la arquitectura hexagonal de puertos y adaptadores propuesta por Cockburn (2024) (\texttt{domain/}, \texttt{application/} e \texttt{infrastructure/}), verificadas mediante reglas de importación en ESLint para que el dominio no dependa de la capa de aplicación ni de la infraestructura, y la aplicación permanezca agnóstica a bibliotecas HTTP o detalles de transporte, tal como se definió en el SAD.

En la \textbf{capa de dominio} (\texttt{sdk/src/domain/}), la entidad \texttt{Transaction} modela el ciclo de vida de un cobro mediante los cinco estados canónicos de \texttt{TransactionStatus} (\texttt{PENDING}, \texttt{APPROVED}, \texttt{DECLINED}, \texttt{VOIDED} y \texttt{ERROR}), expone los predicados de negocio \texttt{isApproved()}, \texttt{isPending()} e \texttt{isFinal()}, conserva en \texttt{rawStatus} la cadena literal devuelta por la pasarela para trazabilidad y auditoría (punto 3) e incluye en \texttt{rejectionReason} la causa normalizada del rechazo (\texttt{RejectionReason}) cuando el estado es \texttt{DECLINED}. Un principio central del diseño es que un rechazo bancario o de tarjeta (\texttt{DECLINED}) constituye un resultado válido del dominio representado como una instancia de \texttt{Transaction}, y nunca se lanza como excepción.

Los objetos de valor garantizan la validez de los datos antes de abrir cualquier conexión de red:

\begin{itemize}
    \item \textbf{\texttt{Amount}}: representa montos monetarios asociados a las cuatro divisas soportadas en \texttt{Currency} (\texttt{COP}, \texttt{USD}, \texttt{EUR} y \texttt{GBP}) conforme al estándar ISO 4217 (International Organization for Standardization [ISO], 2015). En el punto 25 de la bitácora de arquitectura se eliminó la aceptación del tipo numérico \texttt{number} de JavaScript y se restringió la entrada a cadenas decimales canónicas compuestas exclusivamente por dígitos y una parte fraccionaria opcional separada por punto, tanto para evitar errores de representación en aritmética de coma flotante binaria IEEE 754 (IEEE, 2019) como para preservar los ceros decimales finales (\texttt{"15000.00"}), dado que Rapyd calcula su firma HMAC-SHA256 sobre los bytes literales del cuerpo JSON y cualquier alteración del formato decimal invalida la firma (Rapyd Financial Technology, 2024). Los cálculos aritméticos exactos se delegan a \texttt{big-arithmetic.ts} —único archivo del proyecto que importa \texttt{big.js}— y la conversión a unidades menores (centavos para Wompi) se realiza en \texttt{minor-units.ts} desplazando el punto decimal según los exponentes de la norma ISO 4217 (ISO, 2015; punto 33). Adicionalmente, \texttt{Amount} integra un objeto opcional \texttt{TaxBreakdown} (base gravada, IVA e impuesto al consumo), requerido por las pasarelas que exigen desglose tributario en Colombia como Kushki (2024).
    \item \textbf{\texttt{PaymentMethod}}: unión discriminada construida mediante las fábricas estáticas \texttt{PaymentMethod.card()} (que recibe \texttt{token} e \texttt{installments}) y \texttt{PaymentMethod.pse()} (que recibe \texttt{bankCode}, \texttt{userType}, \texttt{docType}, \texttt{docNumber} y \texttt{redirectUrl}). Por diseño de reducción de alcance frente al estándar PCI DSS v4.0 (Payment Card Industry Security Standards Council [PCI SSC], 2022), \texttt{PaymentMethod.card()} acepta exclusivamente un token opaco emitido previamente por la pasarela y jamás recibe el número primario de cuenta (PAN), la fecha de expiración ni el código CVV en el servidor. Para PSE, el campo \texttt{bankCode} acepta tanto el código nativo devuelto por \texttt{getPseBanks()} como el catálogo tipado \texttt{PseBankCode} para pruebas deterministas, cuya traducción hacia los identificadores de Rapyd se concentra en \texttt{rapyd-pse-banks.ts} (punto 68).
    \item \textbf{\texttt{PaymentResult}}: unión discriminada introducida en el punto 39 con las variantes \texttt{TRANSACTION} y \texttt{REDIRECT\_REQUIRED}. Mientras que un cobro con tarjeta en Wompi, Mercado Pago o Kushki se resuelve de forma síncrona devolviendo una \texttt{Transaction}, los cobros por PSE en las cuatro pasarelas y los cobros con tarjeta mediante el formulario alojado (\textit{Hosted Checkout}) de Rapyd requieren que el pagador complete la autenticación en una página externa; en esos casos el SDK retorna un objeto \texttt{PendingRedirect} con el identificador de la orden, la pasarela y la URL de redirección (\texttt{redirectUrl}).
    \item \textbf{\texttt{Credentials} y \texttt{WebhookEvent}}: \texttt{Credentials} valida mediante constructores estáticos (\texttt{forWompi}, \texttt{forMercadoPago}, \texttt{forKushki} y \texttt{forRapyd}) la presencia de las llaves requeridas por cada proveedor y sobrescribe \texttt{toString()} y \texttt{toJSON()} devolviendo \texttt{[REDACTED]} para prevenir fugas accidentales en registros de ejecución. \texttt{WebhookEvent} encapsula la notificación asíncrona verificada (\texttt{gateway}, \texttt{transactionId}, \texttt{status}, \texttt{amount}, \texttt{timestamp}) e incorporó en la versión \texttt{0.7.0-rc.1} los campos opcionales \texttt{rawStatus} y \texttt{orderReference} para permitir que los comercios concilien órdenes pendientes cuando la pasarela notifica un identificador de pago distinto al identificador inicial de sesión o de \textit{checkout} (punto 118).
\end{itemize}

Los servicios de dominio comprenden la verificación de notificaciones entrantes (\texttt{WebhookVerifier} y los cuatro verificadores específicos en \texttt{sdk/src/\allowbreak domain/\allowbreak services/}) y las tablas de traducción de estados nativos (\texttt{native-status.ts}). La comparación de resúmenes criptográficos se ejecuta en \texttt{signature-utils.ts} mediante \texttt{crypto.timingSafeEqual} sobre \textit{buffers} de igual longitud para impedir ataques de temporización lateral, aplicando además una ventana antirreproducción (\textit{anti-replay}) de 300 segundos sobre la marca temporal del evento. En el punto 15 se añadió el parámetro opcional \texttt{gateway} a \texttt{verifyWebhook(payload, headers, secret, gateway?)} para que un comercio que migra de pasarela pueda enrutar la verificación desde su configuración sin alterar la lógica de negocio cuando dos proveedores comparten cabeceras genéricas. Asimismo, en el punto 70 se documentó y reflejó en \texttt{MercadoPagoWebhookHandler} el modelo de conciliación en dos pasos de Mercado Pago (2024): dado que el cuerpo del webhook de Mercado Pago notifica el identificador del recurso pero no transporta el estado financiero definitivo, el verificador valida la firma \texttt{x-signature} y emite el evento con estado \texttt{PENDING} para que el comercio consulte el estado autoritativo mediante \texttt{getPaymentStatus()}.

El catálogo de errores técnicos se unifica en la clase \texttt{KitPagosError}, cuyos trece códigos (\texttt{KitPagosErrorCode}) distinguen fallos de validación de entrada (\texttt{VALIDATION\_ERROR}, \texttt{UNSUPPORTED\_CURRENCY}, \texttt{UNSUPPORTED\_OPERATION}), fallos de autenticación y firma (\texttt{AUTHENTICATION\_ERROR}, \texttt{INVALID\_SIGNATURE}), fallos transitorios de red o capacidad (\texttt{NETWORK\_ERROR}, \texttt{TIMEOUT}, \texttt{RATE\_LIMIT}, \texttt{GATEWAY\_UNAVAILABLE}) y errores de negocio o de contrato del proveedor (\texttt{TRANSACTION\_NOT\_FOUND}, \texttt{DUPLICATE\_PAYMENT}, \texttt{GATEWAY\_ERROR}, \texttt{UNKNOWN\_ERROR}).

En la \textbf{capa de aplicación} (\texttt{sdk/src/application/}), el puerto de salida \texttt{PaymentGatewayPort} define las tres operaciones que todo adaptador debe satisfacer: \texttt{createPayment(request)}, \texttt{getStatus(transactionId)} y \texttt{getPseBanks()}. Siguiendo el patrón de integración \textit{Message Translator} hacia un modelo canónico de datos (Hohpe \& Woolf, 2003), la traducción de las respuestas heterogéneas de cada proveedor hacia las entidades del dominio recae en \texttt{ResponseNormalizer}, que delega en cuatro módulos especializados (\texttt{wompi-normalizer.ts}, \texttt{mercadopago-normalizer.ts}, \texttt{kushki-normalizer.ts} y \texttt{rapyd-normalizer.ts}, separados en el punto 34). La resiliencia ante fallos transitorios se implementa en \texttt{RetryHandler} mediante retroceso exponencial con fluctuación aleatoria (\textit{exponential backoff with jitter}) y una función \texttt{sleep} inyectable para pruebas deterministas; siguiendo la decisión registrada en el punto 35 de la bitácora de arquitectura, la fachada aplica \texttt{RetryHandler} exclusivamente sobre las operaciones idempotentes de lectura (\texttt{getPaymentStatus} y \texttt{getPseBanks}) y jamás reintenta \texttt{createPayment}, dado que repetir automáticamente un \texttt{POST} de creación tras un corte de red o un \texttt{504 Gateway Timeout} podría generar un doble cobro en las pasarelas que no comparten llave de idempotencia entre intentos. Por último, \texttt{ErrorHandler} y su módulo auxiliar \texttt{gateway-error-detail.ts} traducen los códigos de estado HTTP y las estructuras de error nativas a \texttt{KitPagosError}, redactan cualquier credencial presente en la configuración de Axios, tratan el código HTTP \texttt{423} (\texttt{RESOURCE\_LOCKED} en Mercado Pago) como transitorio y reintentable (punto 100) e incorporan desde la versión \texttt{0.6.0-rc.5} el código y motivo nativos de la pasarela en el mensaje de error para facilitar el diagnóstico sin exponer secretos (punto 107).

En la \textbf{capa de infraestructura} (\texttt{sdk/src/infrastructure/}), la fachada \texttt{KitPagos} actúa como único punto de entrada del servidor (Gamma et al., 1994), valida la configuración mediante \texttt{SdkConfigurator}, instancia perezosamente los adaptadores a través de \texttt{GatewayFactory} y coordina el flujo entre validaciones de dominio, ejecución del adaptador, política de reintentos y manejo de errores. En el punto 57, \texttt{SdkConfigurator} se amplió para que la propiedad \texttt{baseUrl} acepte tanto una cadena única (cuando todas las llamadas se dirigen a la API de Simulación) como un mapa parcial por pasarela (\texttt{Partial\textless{}Record\textless{}Gateway, string\textgreater{}\textgreater{}}), permitiendo que un mismo proceso apunte simultáneamente a los cuatro sandboxes oficiales o combine pasarelas reales con simuladas.

\subsubsection{Los cuatro adaptadores de pasarela}

Cada adaptador implementa el patrón \textit{Adapter} (Gamma et al., 1994) para encapsular las particularidades de autenticación, firmado criptográfico, estructura de carga útil y consulta de estado de su proveedor, como se resume en la Tabla \ref{tab:adaptadores-sdk}:

\begin{itemize}
    \item \textbf{\texttt{WompiAdapter}} (\texttt{WompiAdapter.ts} y \texttt{wompi-pse.ts}): autentica las peticiones enviando la llave pública del comercio (\texttt{pub\_test\_*} o \texttt{pub\_prod\_*}) como portador en la cabecera \texttt{Authorization: Bearer} conforme a la especificación de Wompi Colombia (2024), mientras que la llave privada no se transmite por la red en ninguna operación del flujo (punto 50). Antes de cada cobro, el adaptador realiza una petición \texttt{GET /merchants/:publicKey} para obtener un token de aceptación de términos (\texttt{acceptance\_token}) de un solo uso, exigido por Wompi en cada transacción y no reutilizable en caché. A continuación, convierte el monto a centavos enteros (\texttt{amount\_in\_cents}) y calcula la firma de integridad SHA-256 concatenando exactamente \texttt{reference + amount\_in\_cents + currency + integritySecret} (punto 44). En los cobros por PSE, dado que Wompi genera la URL del banco de manera asíncrona tras crear la transacción en estado \texttt{PENDING}, \texttt{wompi-pse.ts} realiza hasta cinco sondeos espaciados sobre \texttt{GET /transactions/:id} hasta que aparece \texttt{payment\_method.extra.async\_payment\_url}; a partir de la versión \texttt{0.4.0}, el normalizador solo clasifica una transacción \texttt{PENDING} como \texttt{REDIRECT\_REQUIRED} cuando dicha URL está efectivamente presente, evitando que los cobros con tarjeta que nacen pendientes se confundan con redirecciones bancarias.
    \item \textbf{\texttt{MercadoPagoAdapter}} (\texttt{MercadoPagoAdapter.ts}, \texttt{mercadopago-card.ts} y \texttt{mercadopago-pse.ts}): autentica mediante el token de acceso en la cabecera \texttt{Authorization: Bearer} y adjunta en cada petición mutante una cabecera \texttt{X-Idempotency-Key} con un identificador UUID v4 generado por llamada (Mercado Pago, 2024; punto 48). A raíz de las mediciones empíricas en el sandbox de Colombia, donde el endpoint clásico \texttt{/v1/payments} rechazaba las credenciales de prueba por políticas internas de evaluación de fraude (\texttt{cc\_rejected\_high\_risk}), en el punto 96 (versión \texttt{0.6.0}) se migró también el flujo de tarjeta a la API de órdenes (\texttt{POST /orders}) utilizando credenciales de un usuario de prueba vendedor (\texttt{APP\_USR-*}, punto 95). Para construir la orden de tarjeta sin pedirle datos redundantes al comercio, \texttt{mercadopago-card.ts} consulta los primeros seis dígitos del BIN en \texttt{GET /v1/card\_tokens/:token} y resuelve el \texttt{payment\_method\_id} (\texttt{visa}, \texttt{master}, \texttt{amex}) en \texttt{GET /v1/payment\_methods/search}. Cuando el procesador rechaza el cobro con tarjeta, Mercado Pago responde con estado HTTP \texttt{402 Payment Required} e incluye el cuerpo completo de la orden con \texttt{status: "failed"} y \texttt{status\_detail}; el adaptador intercepta esa respuesta y la entrega al normalizador para devolver una \texttt{Transaction} en estado \texttt{DECLINED} con su \texttt{rejectionReason} en lugar de lanzar una excepción técnica (puntos 83 y 96). En PSE (\texttt{mercadopago-pse.ts}), el adaptador normaliza las direcciones de bucle local IPv6 (\texttt{::1} y \texttt{::ffff:127.0.0.1}) a \texttt{127.0.0.1} para evitar el error \texttt{invalid\_ip\_address} del validador de Mercado Pago (puntos 96 y 97), extrae la URL bancaria desde \texttt{transactions.payments[0].payment\_method.ticket\_url} y, cuando el sandbox deja una orden PSE en estado \texttt{processing} sin emitir \texttt{ticket\_url}, retorna una \texttt{Transaction} en estado \texttt{PENDING} para que el comercio pueda sondearla posteriormente (punto 100). La consulta de estado \texttt{getStatus()} distingue por el prefijo \texttt{ORD} si debe invocar \texttt{GET /orders/:id} o \texttt{GET /v1/payments/:id}.
    \item \textbf{\texttt{KushkiAdapter}} (\texttt{KushkiAdapter.ts}, \texttt{kushki-charge.ts}, \texttt{kushki-pse.ts} y \texttt{kushki-amount.ts}): separa las cabeceras \texttt{Public-Merchant-Id} (utilizada para consultar el catálogo de bancos en \texttt{GET /transfer/v1/bankList}) y \texttt{Private-Merchant-Id} (utilizada en cobros y consultas de estado en el servidor), de acuerdo con el modelo de credenciales de Kushki (2024). A diferencia de las demás pasarelas, Kushki exige en \texttt{kushki-amount.ts} un objeto estructurado con \texttt{subtotalIva}, \texttt{subtotalIva0}, \texttt{iva} y \texttt{currency}, derivado del \texttt{TaxBreakdown} de \texttt{Amount}. En los cobros con tarjeta (\texttt{POST /card/v1/charges} con \texttt{fullResponse: true}), las pruebas de contrato revelaron dos comportamientos particulares de Kushki en Colombia: cuando el cobro requiere validación asíncrona, la respuesta devuelve \texttt{ticketNumber} pero su estado posterior solo puede consultarse en \texttt{GET /card-async/v1/status/:ticketNumber} y no en \texttt{/charges/} (punto 54); y cuando el procesador bancario declina la tarjeta en el sandbox UAT (tarjetas de prueba terminadas en \texttt{017}, \texttt{019}, \texttt{021}, \texttt{022} y \texttt{023}), Kushki responde con código HTTP \texttt{400 Bad Request} acompañado de un \texttt{details.responseCode} numérico y un \texttt{ticketNumber}, reservando los códigos alfanuméricos \texttt{K*} (\texttt{K005}, \texttt{K322}) para errores técnicos de formato o configuración. Por ello, en el punto 122 (versión \texttt{0.7.0-rc.2}), \texttt{kushki-charge.ts} incorporó el predicado \texttt{acceptDeclinedCharge()} en \texttt{validateStatus}, aceptando los HTTP \texttt{400} con \texttt{responseCode} numérico distinto de \texttt{000} y \texttt{ticketNumber} para que el normalizador produzca una \texttt{Transaction} con estado \texttt{DECLINED} y su \texttt{RejectionReason} correspondiente. En PSE (\texttt{kushki-pse.ts}), dado que el comercio entrega al SDK los datos bancarios y de documento del pagador, el adaptador ejecuta internamente las dos etapas exigidas por Kushki (\texttt{POST /transfer/v1/tokens} seguido de \texttt{POST /transfer/v1/init}, punto 50), mientras que \texttt{getStatus()} recorre secuencialmente \texttt{/transfer/v1/status/}, \texttt{/card-async/v1/status/} y \texttt{/charges/} cuando las rutas previas responden con \texttt{404} o \texttt{400} (puntos 53 y 54).
    \item \textbf{\texttt{RapydAdapter}} (\texttt{RapydAdapter.ts}, \texttt{rapyd-signature.ts}, \texttt{rapyd-checkout.ts}, \texttt{rapyd-pse.ts} y \texttt{rapyd-payload.ts}): firma cada petición HTTP en \texttt{rapyd-signature.ts} generando una sal aleatoria (\texttt{salt}) y una marca temporal Unix (\texttt{timestamp}), concatenando \texttt{httpMethod + urlPath + salt + timestamp + accessKey + secretKey + bodyString}, calculando el HMAC-SHA256 en hexadecimal y codificando esa cadena hexadecimal en Base64 conforme a la guía de seguridad de Rapyd Financial Technology (2024) (puntos 31 y 33). Para evitar discrepancias de espacios o de orden de llaves entre la cadena firmada y la enviada por red, el adaptador serializa el cuerpo con \texttt{JSON.stringify} una única vez y pasa esa misma cadena a la firma y a Axios. En los cobros con tarjeta, dado que Rapyd exige certificación PCI DSS nivel 1 o el uso de su formulario alojado (\textit{Hosted Checkout}) y su librería de navegador requiere crear previamente el \textit{checkout} desde el servidor, \texttt{rapyd-checkout.ts} invoca \texttt{POST /v1/checkout} y devuelve \texttt{REDIRECT\_REQUIRED} con la URL \texttt{redirect\_url} de la página alojada por Rapyd (PCI SSC, 2022; punto 50). En los cobros por PSE (\texttt{rapyd-pse.ts}), el sandbox de Rapyd no dispone de un catálogo de bancos independiente, sino que modela cada entidad financiera de Colombia como un tipo de método de pago distinto bajo el patrón \texttt{co\_pse\_*\_bank} y exige un identificador de cliente \texttt{cus\_*} con documento de identidad colombiano; por tanto, el adaptador filtra \texttt{GET /v1/payment\_methods/countries/co} y, en el momento del cobro, crea primero el cliente en \texttt{POST /v1/customers} y luego emite el pago en \texttt{POST /v1/payments} (puntos 19 y 68). Tras verificar en el punto 102 que de los 47 métodos \texttt{co\_pse\_*\_bank} listados por Rapyd en el sandbox solo nueve aceptan realmente cargos sin devolver \texttt{ERROR\_CREATE\_PAYMENT}, el adaptador filtra en entorno \texttt{sandbox} la respuesta de \texttt{getPseBanks()} contra el conjunto \texttt{RAPYD\_SANDBOX\_PSE\_TYPES}. Finalmente, en la normalización de estados, un pago de Rapyd en estado \texttt{CLO} (\textit{Closed}) solo se traduce a \texttt{APPROVED} cuando el indicador booleano \texttt{paid} es verdadero.
\end{itemize}

\begin{table}[H]
\centering
\small
\caption{Síntesis comparativa de la implementación de los cuatro adaptadores en el SDK}
\label{tab:adaptadores-sdk}
\begin{tabular}{|p{2.3cm}|p{3.3cm}|p{4.1cm}|p{4.1cm}|}
\hline
\textbf{Pasarela} & \textbf{Autenticación e integridad} & \textbf{Flujo de tarjeta} & \textbf{Flujo de PSE y consulta} \\
\hline
Wompi (\texttt{WompiAdapter}) & \texttt{Bearer :publicKey} + token de aceptación único + SHA-256 (\texttt{ref + cents + curr + secret}). & \texttt{POST /transactions} con token \texttt{tok\_*}; respuesta síncrona (\texttt{TRANSACTION}). & \texttt{POST /transactions} + hasta 5 sondeos por \texttt{async\_payment\_url}; estado en \texttt{GET /transactions/:id}. \\
\hline
Mercado Pago (\texttt{MercadoPagoAdapter}) & \texttt{Bearer :accessToken} (\texttt{APP\_USR-*}) + \texttt{X-Idempotency-Key} UUID v4 por llamada. & Consulta de BIN + \texttt{POST /orders}; HTTP \texttt{402} normalizado a \texttt{DECLINED} (punto 96). & \texttt{POST /orders} con IP normalizada y \texttt{ticket\_url}; consulta por prefijo \texttt{ORD} en \texttt{/orders/:id}. \\
\hline
Kushki (\texttt{KushkiAdapter}) & \texttt{Public-Merchant-Id} (bancos) y \texttt{Private-Merchant-Id} (cargos y estado). & \texttt{POST /card/v1/charges} con desglose IVA; HTTP \texttt{400} con código numérico normalizado a \texttt{DECLINED} (punto 122). & Dos pasos (\texttt{/transfer/v1/tokens} + \texttt{init}); consulta en cascada por \texttt{transfer}, \texttt{card-async} y \texttt{charges}. \\
\hline
Rapyd (\texttt{RapydAdapter}) & Cabeceras \texttt{access\_key}, \texttt{salt}, \texttt{timestamp} y \texttt{signature} (HMAC-SHA256 hex en Base64). & \textit{Hosted Checkout} (\texttt{POST /v1/checkout}) que retorna \texttt{REDIRECT\_REQUIRED} (punto 50). & Dos pasos (\texttt{POST /v1/customers} + \texttt{POST /v1/payments} con \texttt{co\_pse\_*\_bank}); filtro de 9 bancos en sandbox (punto 102). \\
\hline
\end{tabular}
\end{table}

\subsubsection{Módulo de tokenización en navegador}

Para que el servidor del comercio reciba únicamente tokens opacos sin entrar en contacto con datos sensibles de tarjeta (PAN y CVV), se construyó el módulo independiente \texttt{sdk/src-browser/}, exportado bajo el subcamino \texttt{kit-pagos-colombia/\allowbreak browser} y empaquetado mediante \texttt{esbuild} (\texttt{sdk/scripts/build-browser.js}) en tres formatos sin dependencias externas: ES Module (\texttt{dist/browser/index.js}), CommonJS (\texttt{dist/browser/index.cjs}) e IIFE autocontenido para etiquetas \texttt{\textless{}script\textgreater{}} (\texttt{dist/browser/\allowbreak kit-pagos-browser.iife.js}, que expone el objeto global \texttt{KitPagosBrowser}) (punto 77).

La clase \texttt{KitPagosBrowser} ofrece una única operación asíncrona, \texttt{tokenizeCard(input, options)}, que recibe los datos de la tarjeta capturados en el navegador, la pasarela destino, la llave pública del comercio y el entorno (\texttt{sandbox}, \texttt{production} o \texttt{simulator}). Por razones de seguridad de frontera, el módulo no permite configurar URLs arbitrarias desde el cliente: resuelve el destino exclusivamente a partir de un catálogo cerrado de dominios oficiales por pasarela y entorno (\texttt{https://sandbox.wompi.co/v1}, \texttt{https://api.mercadopago.com}, \texttt{https://api-uat.kushkipagos.com}) o hacia \texttt{http://localhost:3000/v1/sim/:gateway} cuando el entorno es \texttt{simulator}, impidiendo que una configuración maliciosa redirija el PAN hacia un servidor externo (puntos 77 y 78).

Internamente, \texttt{KitPagosBrowser} delega en tres tokenizadores directos por \texttt{fetch}:

\begin{enumerate}
    \item \textbf{\texttt{WompiTokenizer}}: invoca \texttt{POST /tokens/cards} autenticando con la llave pública en \texttt{Authorization: Bearer} y normalizando el año de expiración a dos dígitos (\texttt{exp\_year}) conforme al contrato de Wompi Colombia (2024) (punto 81).
    \item \textbf{\texttt{MercadoPagoTokenizer}}: tras verificar empíricamente en el punto 78 que \texttt{POST https://api.mercadopago.com/v1/card\_tokens?public\_key=:pk} responde en producción con cabeceras CORS abiertas (\texttt{Access-Control-Allow-Origin: *}) y emite tokens de 32 caracteres sin necesidad de cargar el SDK pesado \texttt{MercadoPago.js}, el tokenizador envía el PAN, el titular y el documento de identificación colombiano (\texttt{docType} y \texttt{docNumber}, obligatorios en Mercado Pago Colombia) directamente por REST (Mercado Pago, 2024).
    \item \textbf{\texttt{KushkiTokenizer}}: incorporado en la versión \texttt{0.6.0-rc.4} (punto 106) después de comprobar con una petición \texttt{OPTIONS} de preflight y un \texttt{POST} real desde el origen \texttt{http://localhost:5174} que \texttt{POST https://api-uat.kushkipagos.com/card/v1/tokens} incluye \texttt{Access-Control-Allow-Origin: *} y acepta la cabecera \texttt{Public-Merchant-Id} (Kushki, 2024). Dado que Kushki vincula cada token de tarjeta al monto y moneda de la compra en el momento de emitirlo, \texttt{KushkiTokenizer} exige los campos \texttt{amount} y \texttt{currency} en la entrada de \texttt{tokenizeCard()}.
\end{enumerate}

Cuando se invoca \texttt{tokenizeCard()} para \texttt{Gateway.RAPYD}, el módulo lanza inmediatamente un \texttt{KitPagosBrowserError} con código \texttt{UNSUPPORTED\_OPERATION}, dado que Rapyd opera mediante \textit{Hosted Checkout} iniciado desde el servidor (Rapyd Financial Technology, 2024; puntos 50 y 106). En el punto 114 de la bitácora de arquitectura se dejó explícito el límite normativo de este enfoque frente a PCI DSS v4.0 (PCI SSC, 2022): al capturar el PAN en campos de formulario controlados por el comercio antes de enviarlo directamente a la pasarela, el comercio se ubica bajo el cuestionario de autoevaluación SAQ A-EP (y no SAQ A, que requeriría \textit{iframes} alojados por cada proveedor) y es responsable de declarar las directivas \texttt{Content-Security-Policy} (\texttt{connect-src}) en las respuestas HTML de su servidor.

\subsubsection{Versiones publicadas en npm}

El paquete \texttt{kit-pagos-colombia} se publicó en el registro oficial de npm (\texttt{https://www.npmjs.com/package/kit-pagos-colombia}) siguiendo versionado semántico bajo la serie \texttt{0.y.z} durante la etapa experimental de la tesis. En la Tabla \ref{tab:versiones-npm} se detalla la evolución de las versiones publicadas y los hitos técnicos que motivaron cada corte.

\begin{table}[H]
\centering
\small
\caption{Evolución de las versiones del paquete \texttt{kit-pagos-colombia} publicadas en npm}
\label{tab:versiones-npm}
\begin{tabular}{|p{2.2cm}|p{2.0cm}|p{9.6cm}|}
\hline
\textbf{Versión} & \textbf{Canal npm} & \textbf{Cambios técnicos principales incorporados} \\
\hline
\texttt{0.1.0} & \texttt{latest} & Primera publicación pública del núcleo hexagonal tras rebajar la versión inicial desde \texttt{1.0.0}; incluye \texttt{Transaction}, \texttt{Amount}, los cuatro adaptadores base y \texttt{WebhookVerifier}. \\
\hline
\texttt{0.2.0} & \texttt{latest} & Incorporación del punto de entrada \texttt{kit-pagos-colombia/\allowbreak browser} en \texttt{exports} de \texttt{package.json} y del script de verificación prepublicación \texttt{check-published-package.ts} (punto 77). \\
\hline
\texttt{0.3.0} -- \texttt{0.3.1} & \texttt{latest} & Soporte para mapa de \texttt{baseUrl} por pasarela en \texttt{SdkConfigurator} (punto 57), catálogo tipado \texttt{PseBankCode} con traducción para Rapyd (punto 68) y conciliación de webhooks de Mercado Pago en estado \texttt{PENDING} (punto 70). \\
\hline
\texttt{0.4.0} & \texttt{latest} & Corrección en \texttt{WompiAdapter} para devolver \texttt{REDIRECT\_REQUIRED} únicamente cuando la transacción \texttt{PENDING} contiene \texttt{async\_payment\_url}, evitando falsas redirecciones en cobros con tarjeta. \\
\hline
\texttt{0.5.0} & \texttt{latest} & Inclusión del flujo de trabajo automatizado \texttt{.github/\allowbreak workflows/\allowbreak publish.yml} con procedencia (\texttt{--provenance}) y verificación de compilación de fragmentos de código (\texttt{check:readme}, punto 76). \\
\hline
\texttt{0.6.0} (\texttt{rc.1}--\texttt{rc.5}) & \texttt{next} $\rightarrow$ \texttt{latest} & Migración de tarjeta en Mercado Pago a \texttt{POST /orders} con normalización de IPv6 y manejo de HTTP \texttt{402} como \texttt{DECLINED} (punto 96); manejo de PSE en \texttt{processing} sin URL y HTTP \texttt{423} reintentable (punto 100); filtro de los 9 bancos PSE cobrables en el sandbox de Rapyd (punto 102); incorporación de \texttt{KushkiTokenizer} en \texttt{kit-pagos-colombia/\allowbreak browser} (punto 106); detalle nativo de errores en \texttt{ErrorHandler} (punto 107); corte estable \texttt{0.6.0} en \texttt{latest} (punto 108). \\
\hline
\texttt{0.7.0-rc.1} -- \texttt{0.7.0-rc.2} & \texttt{next} & Adición de \texttt{rawStatus} y \texttt{orderReference} en \texttt{WebhookEvent} y soporte de \texttt{PAYMENT\_UPDATED} en Rapyd para conciliación de webhooks reales (punto 118); aceptación de HTTP \texttt{400} con \texttt{responseCode} numérico en \texttt{KushkiAdapter} como transacción \texttt{DECLINED} en lugar de excepción técnica (punto 122). \\
\hline
\end{tabular}
\end{table}

\subsection{API de Simulación}

\subsubsection{Arquitectura del servidor y separación de superficies: \texttt{/v1/sim} y \texttt{/v1/api}}

La API de Simulación (\texttt{simulator-api/}) se construyó sobre Fastify 4 y TypeScript aplicando los principios del estilo arquitectónico REST (Fielding, 2000) y separando la fábrica de la aplicación en \texttt{src/app.ts} (\texttt{buildApp()}) del arranque del socket en \texttt{src/server.ts} para permitir pruebas de integración rápidas en memoria mediante \texttt{app.inject()} sin abrir puertos TCP (punto 10). El proceso expone dos superficies REST claramente diferenciadas que atienden a los dos patrones de integración del proyecto definidos en el SAD (punto 125):

La primera superficie, \textbf{\texttt{/v1/sim/:gateway/*}} (\texttt{simulator-api/src/routes/}), implementa el patrón de integración P1 (\textit{SDK directo}) actuando como un doble de pruebas determinista (Ammann \& Offutt, 2016) que reproduce el contrato HTTP nativo de cada proveedor (Kushki, 2024; Mercado Pago, 2024; Rapyd Financial Technology, 2024; Wompi Colombia, 2024). Comprende 33 rutas de pasarela —siete de Wompi (\texttt{wompi.ts}), nueve de Mercado Pago (\texttt{mercadopago.ts} y \texttt{mercadopago-pse.ts}), nueve de Rapyd (\texttt{rapyd.ts} y \texttt{rapyd-Flow.ts}) y ocho de Kushki (\texttt{kushki.ts})— más el disparador de notificaciones \texttt{POST /v1/sim/webhooks/trigger} (\texttt{webhooks.ts}) y la ruta de salud \texttt{GET /health}, sumando 35 endpoints. Cuando el SDK apunta su \texttt{baseUrl} a \texttt{/v1/sim/:gateway}, ejecuta exactamente el mismo código de firmado, serialización, cabeceras y normalización que cuando llama a los sandboxes reales. Además de los endpoints JSON, esta superficie sirve páginas HTML interactivas (\texttt{src/routes/helpers/bankPage.ts}) para simular tanto el portal bancario PSE de las cuatro pasarelas (permitiendo aprobar o rechazar la transferencia en el navegador y redirigir de vuelta a la tienda del comercio, punto 93) como la página alojada de tarjeta de Rapyd en \texttt{/v1/sim/rapyd/checkout/:id/pay} (puntos 92 y 106).

El almacenamiento de \texttt{/v1/sim} reside en memoria dentro de \texttt{GatewayStores} (\texttt{src/storage/GatewayStores.ts}), que mantiene ocho colecciones tipadas (\texttt{InMemoryStore\textless{}T\textgreater{}}) separadas por pasarela y tipo de recurso (transacciones de Wompi, pagos y órdenes de Mercado Pago, pagos y \textit{checkouts} de Rapyd, y cargos, tokens e inicializaciones de transferencia de Kushki), junto con tres mapas auxiliares: \texttt{CardTokenOutcomes} (que preserva el escenario determinado al tokenizar una tarjeta en el navegador para aplicarlo cuando el servidor ejecuta el cargo, punto 83), \texttt{IdempotencyKeys} (que devuelve la respuesta original cuando Mercado Pago recibe la misma cabecera \texttt{X-Idempotency-Key} dos veces) y \texttt{BankReturnUrls} (que conserva el \texttt{returnUrl} del comercio durante el flujo de redirección PSE). Las transiciones de ciclo de vida de cada entidad están gobernadas por máquinas de estados declarativas (\texttt{src/state/TransitionEngine.ts} y las tablas \texttt{wompiMachine.ts}, \texttt{mercadopagoMachine.ts}, \texttt{rapydMachine.ts} y \texttt{kushkiMachine.ts}), las cuales lanzan \texttt{InvalidTransitionError} ante cualquier intento de mutar una transacción que ya alcanzó un estado terminal (punto 8).

La segunda superficie, \textbf{\texttt{/v1/api/*}} (\texttt{simulator-api/src/kit-pagos-api/}), materializa el patrón de integración P2 (\textit{API REST centralizada}) ofreciendo cinco rutas unificadas para comercios cuyo backend no está escrito en TypeScript o que prefieren consumir un contrato JSON único:

\begin{enumerate}
    \item \textbf{\texttt{GET /v1/api/gateways}}: expone el catálogo de pasarelas disponibles, el entorno activo por proveedor (\texttt{simulator}, \texttt{sandbox} o \texttt{production}), si requiere credenciales en la cabecera de la petición (\texttt{requiresCredentialsHeader}) y el modo de captura de tarjeta (\texttt{cardCapture: "BROWSER\_TOKEN"} para Wompi, Mercado Pago y Kushki, frente a \texttt{"HOSTED\_PAGE"} para Rapyd, añadido en el punto 106).
    \item \textbf{\texttt{POST /v1/api/payments}}: recibe un cuerpo unificado validado con Zod (\texttt{gateway}, \texttt{amount}, \texttt{currency}, \texttt{paymentMethod}, \texttt{reference}, \texttt{customer}, \texttt{tax}), extrae y normaliza la dirección IP del comprador en \texttt{payer-ip.ts} cuando el pago es PSE y falta \texttt{payerIp} (respetando \texttt{X-Forwarded-For} únicamente cuando \texttt{SIMULATOR\_TRUST\_FORWARDED\_FOR=true}, puntos 92 y 97), e invoca \texttt{kitPagos.createPayment()}.
    \item \textbf{\texttt{GET /v1/api/payments/:id}}: consulta el estado normalizado de una transacción en la pasarela indicada por el parámetro de consulta \texttt{gateway}.
    \item \textbf{\texttt{GET /v1/api/pse-banks}}: devuelve la lista homologada de bancos PSE disponibles para la pasarela solicitada.
    \item \textbf{\texttt{POST /v1/api/webhooks/:gateway}}: registra un analizador de cuerpo en crudo (\textit{raw body parser}) para preservar los bytes exactos del JSON recibido, delega la validación criptográfica en \texttt{kitPagos.verifyWebhook()} y retorna el \texttt{WebhookEvent} normalizado (punto 63).
\end{enumerate}

La resolución de credenciales y del entorno destino en \texttt{/v1/api} sigue el modelo híbrido implementado en \texttt{CredentialResolver.ts}, \texttt{targetEnvironment.ts} y \texttt{KitPagosProvider.ts} (puntos 63, 66, 69, 80 y 107). El operador o comercio puede suministrar credenciales por petición mediante la cabecera \texttt{X-Gateway-Credentials} (codificada en JSON o Base64) o delegar en las variables de entorno del servidor declaradas en la tabla \texttt{SERVER\_PROFILES}. La función \texttt{resolveTargetEnvironment()} aplica cuatro reglas deterministas: el valor \texttt{production} en \texttt{X-Target-Environment} nunca es sobreescrito por variables del servidor y exige credenciales explícitas en la petición; en \texttt{simulator}, si no se envían credenciales, se inyectan credenciales ficticias (\texttt{MOCK\_CREDENTIALS}) y el SDK apunta internamente a \texttt{SIMULATOR\_SDK\_BASE\_URL} (\texttt{http://127.0.0.1::port/v1/sim}); y en \texttt{sandbox}, el SDK invoca los dominios reales de prueba de cada proveedor. Cuando la petición utiliza el perfil del servidor sin cabeceras por llamada, \texttt{KitPagosProvider} reutiliza una instancia \textit{singleton} (Gamma et al., 1994) de \texttt{KitPagos}; cuando recibe credenciales o escenarios por petición, construye una instancia efímera aislada para evitar contaminación de estado entre clientes concurrentes.

Para proteger las credenciales que transitan tanto hacia \texttt{/v1/api} como hacia \texttt{/v1/sim}, en el punto 120 de la bitácora de arquitectura se reemplazó la lista negra de cabeceras del registrador Pino por una lista blanca estricta (\texttt{LOGGABLE\_HEADERS}) en \texttt{src/utils/redactSerializer.ts}: el servidor solo conserva en claro diez cabeceras operativas inocuas (\texttt{host}, \texttt{user-agent}, \texttt{content-type}, \texttt{content-length}, \texttt{accept}, \texttt{accept-encoding}, \texttt{accept-language}, \texttt{origin}, \texttt{referer} y \texttt{x-request-id}) y sustituye cualquier otra cabecera —incluidas \texttt{authorization}, \texttt{x-gateway-credentials}, \texttt{private-merchant-id}, \texttt{public-merchant-id}, \texttt{access\_key} y \texttt{signature}— por la cadena literal \texttt{[REDACTED]}. Asimismo, en el punto 125 se integró \texttt{@fastify/swagger} y \texttt{@fastify/swagger-ui} para servir en \texttt{/docs} la especificación OpenAPI 3.0.3 interactiva de \texttt{/v1/api} y de \texttt{/v1/sim}, ocultando de la interfaz pública las rutas marcadas con \texttt{hide: true}.

\subsubsection{Escenarios deterministas, emisión de webhooks y límites de fidelidad}

Para que las pruebas automatizadas y las demostraciones reproduzcan tanto casos de éxito como rechazos bancarios y fallos técnicos mediante partición del espacio de entrada (Ammann \& Offutt, 2016) sin depender del azar, la función \texttt{scenarioFromRequest()} (\texttt{src/routes/helpers/scenarioFromRequest.ts}) y el detector \texttt{invalidCredential.ts} evalúan cada petición entrante en \texttt{/v1/sim} siguiendo un orden estricto de cinco niveles de precedencia (punto 83):

\begin{enumerate}
    \item \textbf{Cabecera explícita}: si la petición incluye \texttt{x-simulator-scenario} o \texttt{x-simulate-scenario}, se aplica directamente ese escenario (\texttt{APPROVED}, \texttt{DECLINED}, \texttt{PENDING}, \texttt{EXPIRED}, \texttt{VOIDED}, \texttt{TIMEOUT}, \texttt{RATE\_LIMIT}, \texttt{GATEWAY\_ERROR}, \texttt{GATEWAY\_UNAVAILABLE}, \texttt{AUTHENTICATION\_ERROR}, \texttt{DUPLICATE\_PAYMENT}, \texttt{INVALID\_CARD\_TOKEN}, entre otros).
    \item \textbf{Marcador en credencial}: si la llave de autenticación contiene el prefijo o sufijo \texttt{invalid}, el escenario se fija en \texttt{AUTHENTICATION\_ERROR} (\texttt{401}); en Kushki, los marcadores \texttt{inexistente} y \texttt{not\_found} fuerzan \texttt{MERCHANT\_NOT\_FOUND}.
    \item \textbf{Datos de prueba nativos de la pasarela}: reproduce las convenciones documentadas por cada proveedor en sus sandboxes (Kushki, 2024; Mercado Pago, 2024; Wompi Colombia, 2024): en Wompi, tarjetas terminadas en \texttt{4111} declinan y los bancos PSE \texttt{"1"}, \texttt{"2"} y \texttt{"3"} aprueban, declinan o quedan pendientes; en Mercado Pago, los prefijos del nombre del titular \texttt{APRO}, \texttt{CONT}, \texttt{OTHE}, \texttt{FUND}, \texttt{SECU}, \texttt{EXPI} y \texttt{FORM} determinan el resultado; en Kushki, las tarjetas terminadas en \texttt{017}, \texttt{019}, \texttt{021}, \texttt{022} y \texttt{023} devuelven HTTP \texttt{400} con el \texttt{responseCode} numérico de rechazo del procesador (punto 122).
    \item \textbf{Montos reservados en COP}: cuando no concurren las señales anteriores, el monto entero en pesos colombianos actúa como selector determinista de escenario, tal como se sintetiza en la Tabla \ref{tab:montos-simulador}. Adicionalmente, los montos de un pago ya almacenado (\texttt{10600}, \texttt{10602} y \texttt{10604}) fuerzan fallos HTTP \texttt{500}, \texttt{502} y \texttt{504} durante la consulta posterior de estado (\texttt{getStatus}) para ejercitar la política de reintentos de \texttt{RetryHandler}.
    \item \textbf{Valor por omisión}: cualquier otra combinación válida se resuelve como \texttt{APPROVED}.
\end{enumerate}

\begin{table}[H]
\centering
\small
\caption{Catálogo de montos reservados en COP para la selección determinista de escenarios en \texttt{/v1/sim}}
\label{tab:montos-simulador}
\begin{tabular}{|p{3.0cm}|p{4.2cm}|p{6.6cm}|}
\hline
\textbf{Monto (COP)} & \textbf{Escenario activado} & \textbf{Comportamiento en la API de Simulación} \\
\hline
\texttt{10100} & \texttt{DECLINED} & Transacción creada con estado rechazado y motivo \texttt{INSUFFICIENT\_FUNDS} o equivalente nativo. \\
\hline
\texttt{10101} & \texttt{PENDING} & Transacción retenida en estado pendiente de confirmación asíncrona. \\
\hline
\texttt{10102} & \texttt{EXPIRED} & Transacción u orden que transita al estado de expiración por tiempo. \\
\hline
\texttt{10001} -- \texttt{10006} & Motivos específicos de rechazo & Activan \texttt{INSUFFICIENT\_FUNDS}, \texttt{INVALID\_CARD}, \texttt{EXPIRED\_CARD}, \texttt{FRAUD\_SUSPECTED}, \texttt{DO\_NOT\_HONOR} y \texttt{BANK\_REJECTED}. \\
\hline
\texttt{10409} & \texttt{DUPLICATE\_PAYMENT} & Respuesta HTTP \texttt{409 Conflict} por referencia o pago duplicado. \\
\hline
\texttt{10429} & \texttt{RATE\_LIMIT} & Respuesta HTTP \texttt{429 Too Many Requests} sin mutar el almacén. \\
\hline
\texttt{10500}, \texttt{10502}, \texttt{10503}, \texttt{10504} & Fallos de servidor y red & Respuestas HTTP \texttt{500}, \texttt{502}, \texttt{503} y \texttt{504} inmediatas sin crear ni alterar registros. \\
\hline
\texttt{10600}, \texttt{10602}, \texttt{10604} & Fallos en consulta (\texttt{getStatus}) & El pago se crea normalmente, pero las llamadas posteriores de consulta devuelven HTTP \texttt{500}, \texttt{502} o \texttt{504}. \\
\hline
\end{tabular}
\end{table}

El escenario resuelto durante la creación queda fijado en la entidad almacenada (\texttt{scenarioTarget.ts}), de modo que las consultas posteriores o el cierre del flujo de redirección respetan el destino original sin necesidad de reenviar cabeceras. Cuando se solicita una combinación de pasarela, método y escenario que el equipo no pudo medir empíricamente en el sandbox real del proveedor, la tabla de respuestas lanza \texttt{UnsupportedScenarioError} y el servidor responde con HTTP \texttt{501 Not Implemented} en lugar de inventar una carga útil no verificada (punto 83).

La emisión de webhooks soporta dos modalidades complementarias (\texttt{src/webhooks/WebhookEmitter.ts} y \texttt{autoEmission.ts}). En el modo manual, el cliente invoca \texttt{POST /v1/sim/webhooks/trigger} indicando la pasarela, el identificador del pago, el nuevo estado y la URL destino (\texttt{targetUrl}); el emisor aplica la transición en la máquina de estados correspondiente, construye el cuerpo nativo de la notificación, calcula la firma criptográfica exacta que valida \texttt{WebhookVerifier} (SHA-256 con propiedades concatenadas para Wompi, plantilla \texttt{id:\dots;request-id:\dots;ts:\dots;} en HMAC-SHA256 para Mercado Pago, HMAC-SHA256 sobre el cuerpo para Kushki y cadena completa codificada en Base64 para Rapyd) y despacha la petición HTTP con hasta tres reintentos. En el modo automático, activado cuando \texttt{SIMULATOR\_WEBHOOK\_AUTO=true} y \texttt{SIMULATOR\_WEBHOOK\_TARGET\_URL} están definidos, los almacenes registran un gancho \texttt{onStatusChange} en \texttt{attachAutoWebhookEmission()} que dispara la notificación tras cada cambio de estado; para reproducir fielmente el comportamiento observado en el sandbox de Rapyd durante los ensayos en vivo, cuando un pago de Rapyd pasa a aprobado el emisor automático envía en secuencia los tres eventos \texttt{PAYMENT\_SUCCEEDED}, \texttt{PAYMENT\_UPDATED} y \texttt{PAYMENT\_COMPLETED} (puntos 83 y 118).

De acuerdo con el principio de transparencia metodológica de DSR (Peffers et al., 2007), el diseño de la API de Simulación declara explícitamente seis límites de fidelidad frente a los sandboxes de las pasarelas (puntos 48, 50 y 86): (1) el almacenamiento reside en memoria volátil y se reinicia con el proceso; (2) no se ejecutan validaciones criptográficas profundas de emisor bancario ni flujos reales de autenticación 3DS; (3) no se aplican cuotas de tráfico reales salvo cuando el cliente solicita explícitamente el escenario \texttt{RATE\_LIMIT}; (4) las respuestas se producen con latencia local de pocos milisegundos en lugar de los tiempos variables de red de cada proveedor; (5) la página bancaria simulada en \texttt{/v1/sim/*/pse/bank/:id} reemplaza los portales heterogéneos de ACH Colombia por una interfaz HTML controlada; y (6) las combinaciones no medidas en el sandbox real se rechazan con HTTP \texttt{501} en lugar de aproximarse mediante suposiciones.

\subsubsection{Despliegue público en Render y colección de Postman}

El despliegue público de la API de Simulación se definió como infraestructura como código en el archivo \texttt{render.yaml} de la raíz del repositorio y se sirve en \texttt{https://kit-pagos-colombia.onrender.com} a partir de la rama \texttt{devops} (puntos 79, 84, 88 y 122). El servicio compila tanto \texttt{sdk/} como \texttt{simulator-api/} durante la fase de construcción y arranca con \texttt{npm --prefix simulator-api start}, exponiendo \texttt{GET /health} como sonda de disponibilidad. El módulo \texttt{loadServerEnv()} (\texttt{simulator-api/src/config/serverEnv.ts}) carga al inicio el archivo de secretos montado por Render en \texttt{/etc/secrets/.env} sin pisar las variables ya presentes en \texttt{process.env}, configurando \texttt{SIMULATOR\_PUBLIC\_ORIGIN=https://kit-pagos-colombia.onrender.com} (para que los enlaces de redirección PSE y de \textit{Hosted Checkout} devueltos al cliente apunten al dominio público HTTPS), \texttt{SIMULATOR\_SDK\_BASE\_URL=http://127.0.0.1:10000/v1/sim} (para que las llamadas internas del SDK desde \texttt{/v1/api} en modo \texttt{simulator} se resuelvan por bucle local sin salir a internet) y \texttt{SIMULATOR\_TRUST\_FORWARDED\_FOR=true} detrás del proxy inverso de Render. Para permitir que los evaluadores académicos y el prototipo estático consuman el servicio sin aprovisionamiento previo de llaves, el despliegue público opera en modo abierto (sin \texttt{API\_AUTH\_TOKEN}) y resuelve las credenciales de sandbox desde el perfil interno del servidor cuando la petición no incluye \texttt{X-Gateway-Credentials} (punto 88).

Para verificar de extremo a extremo tanto el servidor local como la instancia desplegada en Render, se construyó la colección versionada \texttt{simulator-api/\allowbreak postman/\allowbreak kit-pagos-simulator.postman\_collection.json} junto con los entornos \texttt{local.postman\_environment.json} y \texttt{render.postman\_environment.json} (puntos 88, 106 y 122). La colección agrupa 42 peticiones HTTP y 47 aserciones automatizadas distribuidas en siete carpetas: (1) salud y catálogo de pasarelas (\texttt{GET /health} y \texttt{GET /v1/api/gateways}); (2) los cinco endpoints de la fachada unificada \texttt{/v1/api} con casos de aprobación, rechazo por monto reservado, consulta de estado, listado de bancos PSE y verificación de webhook firmado; (3) cinco llamadas directas a \texttt{/v1/sim} para cada una de las cuatro pasarelas (Wompi, Mercado Pago, Kushki y Rapyd), que incluyen scripts de pre-petición en JavaScript para calcular en tiempo de ejecución la firma SHA-256 de Wompi y la firma HMAC-SHA256 en Base64 de Rapyd, encadenando mediante variables de colección los identificadores dinámicos (\texttt{wompi\_tx\_id}, \texttt{mp\_order\_id}, \texttt{kushki\_ticket}, \texttt{rapyd\_payment\_id}); y (4) una carpeta opcional de humo contra los sandboxes reales a través de \texttt{/v1/api} con \texttt{X-Target-Environment: sandbox}. La colección completa se ejecuta sin fallos mediante Newman (\texttt{npx newman run}) tanto contra \texttt{http://localhost:3000} como contra \texttt{https://kit-pagos-colombia.onrender.com}.

\subsection{Documentación y datos de prueba}

El tercer componente del artefacto es el catálogo estructurado de datos de prueba y evidencia empírica de los sandboxes reales, alojado en el directorio \texttt{docs/testing-data/} y organizado en cuatro módulos de consulta por proveedor junto con las bitácoras de medición en vivo ejecutadas el 18 y el 27 de septiembre de 2026. Su propósito fue eliminar la dispersión de las fuentes oficiales (Kushki, 2024; Mercado Pago, 2024; Rapyd Financial Technology, 2024; Wompi Colombia, 2024) detectada durante la Fase 1 y dotar tanto a los adaptadores del SDK como a las tablas de escenarios de \texttt{/v1/sim} de una fuente de verdad trazable.

Para distinguir con rigor qué comportamientos provienen de mediciones directas del equipo y cuáles provienen únicamente de manuales de los proveedores, el catálogo de \texttt{docs/testing-data/} clasifica cada dato bajo tres niveles de evidencia explícitos:

\begin{itemize}
    \item \textbf{Nivel 1 (Medido en sandbox)}: petición ejecutada contra el sandbox real de la pasarela, con fecha de medición, código de estado HTTP observado, estructura de respuesta capturada y prueba de contrato reproducible en \texttt{sdk/test/sandbox/}.
    \item \textbf{Nivel 2 (Medido hasta el punto de redirección)}: flujo ejecutado en el sandbox real hasta la emisión de la URL de pago externo (\texttt{async\_payment\_url}, \texttt{ticket\_url} o \texttt{redirect\_url}), verificando la creación y la consulta del estado inicial antes de la interacción humana en el navegador.
    \item \textbf{Nivel 3 (Documentado por el proveedor, no verificado en vivo)}: valores de referencia tomados de la documentación oficial de la pasarela que no pudieron ejercitarse programáticamente por restricciones del entorno de pruebas del proveedor.
\end{itemize}

Las campañas de medición registradas en este catálogo revelaron discrepancias sustanciales entre la documentación pública de las pasarelas y el comportamiento real de sus sandboxes en Colombia, hallazgos que determinaron decisiones clave de implementación en el SDK y en la API de Simulación (Tabla \ref{tab:hallazgos-sandboxes}):

\begin{table}[H]
\centering
\small
\caption{Hallazgos empíricos en los sandboxes de las pasarelas y su impacto en el artefacto}
\label{tab:hallazgos-sandboxes}
\begin{tabular}{|p{2.3cm}|p{5.5cm}|p{6.0cm}|}
\hline
\textbf{Pasarela} & \textbf{Comportamiento medido en el sandbox real} & \textbf{Ajuste aplicado en el SDK y en \texttt{/v1/sim}} \\
\hline
Wompi (Wompi Colombia, 2024) & En PSE, cuando se usa el banco \texttt{"1"}, el sandbox entrega la \texttt{async\_payment\_url} y resuelve la transacción a \texttt{APPROVED} en el segundo sondeo sin esperar a que un usuario abra el enlace (punto 43). & \texttt{WompiAdapter} prioriza retornar \texttt{REDIRECT\_REQUIRED} en cuanto aparece \texttt{async\_payment\_url}, y la prueba de contrato acepta tanto \texttt{REDIRECT\_REQUIRED} como \texttt{TRANSACTION} si el sandbox resuelve antes. \\
\hline
Mercado Pago (Mercado Pago, 2024) & Las credenciales de cuenta real en modo prueba son rechazadas por \texttt{/v1/payments} con \texttt{cc\_rejected\_high\_risk} y bloqueadas en \texttt{/v1/orders} por cruce de entornos; además, el conector PSE del sandbox presenta ventanas de intermitencia donde las órdenes nacen en \texttt{processing} sin \texttt{ticket\_url} o devuelven \texttt{423 RESOURCE\_LOCKED} (puntos 45, 95, 100 y 113). & Se aprovisionó un usuario de prueba vendedor (\texttt{APP\_USR-*}) cuyos tokens operan en \texttt{POST /v1/orders} para tarjeta y PSE; el adaptador trata el estado \texttt{processing} sin URL como \texttt{PENDING} y el código HTTP \texttt{423} como reintentable. \\
\hline
Kushki (Kushki, 2024) & Los cargos de tarjeta que requieren validación asíncrona no existen en \texttt{GET /card/v1/charges/:id} y solo responden en \texttt{GET /card-async/v1/status/:ticket}; los rechazos de procesador en tarjeta devuelven HTTP \texttt{400} con \texttt{details.responseCode} numérico (\texttt{017}--\texttt{023}) en lugar de HTTP \texttt{201} (puntos 53, 54 y 122). & \texttt{KushkiAdapter.getStatus()} consulta en cascada \texttt{transfer}, \texttt{card-async} y \texttt{charges}; \texttt{kushki-charge.ts} acepta HTTP \texttt{400} con código numérico y \texttt{ticketNumber} para normalizarlo a \texttt{DECLINED}. \\
\hline
Rapyd (Rapyd Financial Technology, 2024) & El catálogo \texttt{GET /v1/payment\_methods/countries/co} lista 47 métodos \texttt{co\_pse\_*\_bank}, pero 38 de ellos fallan al crear el pago en sandbox con \texttt{ERROR\_CREATE\_PAYMENT}; además, al aprobar un pago, Rapyd emite tres webhooks consecutivos (\texttt{PAYMENT\_SUCCEEDED}, \texttt{PAYMENT\_UPDATED} y \texttt{PAYMENT\_COMPLETED}) donde el estado definitivo viaja en los dos últimos (puntos 19, 102 y 118). & \texttt{RapydAdapter.getPseBanks()} filtra la lista en modo \texttt{sandbox} a los 9 bancos operativos (\texttt{RAPYD\_SANDBOX\_PSE\_TYPES}); \texttt{RapydWebhookHandler} y \texttt{autoEmission.ts} procesan y emiten la secuencia de tres eventos e incluyen \texttt{orderReference}. \\
\hline
\end{tabular}
\end{table}

\subsection{Ejemplos y proyectos prototípicos}

\subsubsection{Ejemplos ejecutables del SDK}

El directorio \texttt{examples/} contiene doce scripts ejecutables en TypeScript que consumen el paquete local mediante \texttt{"kit-pagos-colombia": "file:../sdk"} e importan exclusivamente desde la raíz pública del paquete, sirviendo tanto de guía práctica para desarrolladores como de verificación de contrato en el trabajo \texttt{examples} de la integración continua (puntos 21 y 83).

Los cuatro primeros scripts (\texttt{wompi-card.ts}, \texttt{mercadopago-card.ts}, \texttt{kushki-card.ts} y \texttt{rapyd-card.ts}) muestran el flujo completo de pago con tarjeta en cada pasarela, incluyendo la creación del cobro, la inspección de \texttt{PaymentResult} y la consulta posterior de estado con \texttt{getPaymentStatus()}. Los cuatro siguientes (\texttt{wompi-pse.ts}, \texttt{mercadopago-pse.ts}, \texttt{kushki-pse.ts} y \texttt{rapyd-pse.ts}) recorren el flujo de transferencia bancaria PSE y el manejo del resultado \texttt{REDIRECT\_REQUIRED}. El script \texttt{pse-banks.ts} consulta en paralelo \texttt{getPseBanks()} sobre los cuatro proveedores y muestra la estructura homologada \texttt{PseBank[]}.

Por su parte, \texttt{gateway-interchangeability.ts} demuestra en código la propiedad de intercambiabilidad del patrón P1: una única función \texttt{procesarCobro(sdk, gateway, token)} ejecuta el mismo pago sucesivamente en Wompi, Mercado Pago, Kushki y Rapyd cambiando únicamente el valor del enumerado \texttt{Gateway}, sin una sola sentencia condicional (\texttt{if} o \texttt{switch}) por proveedor. El script interactivo \texttt{demo-live.ts} levanta la API de Simulación en memoria e intercepta el transporte HTTP para imprimir en consola qué cabeceras, firmas criptográficas y cuerpos JSON viajan por el cable en cada llamada. Finalmente, \texttt{simulate-scenarios.ts} ejecuta de forma automatizada una matriz de 61 casos (aprobaciones, rechazos por monto y por datos nativos, redirecciones PSE, errores de autenticación, límites de tasa, caídas de pasarela y verificación de webhooks con firmas válidas e inválidas) contra una instancia en memoria de \texttt{buildApp()}, verificando con código de salida cero que el SDK y el simulador convergen en los 61 escenarios (punto 83).

\subsubsection{Los cuatro prototipos de comercio electrónico y su papel metodológico}

En el directorio \texttt{prototypes/} se construyeron cuatro aplicaciones completas de comercio electrónico en TypeScript que comparten el mismo dominio funcional —un catálogo de seis productos colombianos con precios en COP, selección de pasarela (Wompi, Mercado Pago, Kushki y Rapyd), cobro con tarjeta y con PSE, recepción de webhooks firmados, página de retorno bancaria y panel de trazabilidad HTTP en vivo (\texttt{trace.ts}) que registra cada petición y respuesta cruzando la frontera del comercio—. En el punto 124 de la bitácora de arquitectura se consolidó la nomenclatura definitiva de los cuatro prototipos en cuatro carpetas autoexplicativas (Tabla \ref{tab:prototipos-resumen}):

\begin{table}[H]
\centering
\small
\caption{Los cuatro prototipos implementados en \texttt{prototypes/} y su papel en la investigación}
\label{tab:prototipos-resumen}
\begin{tabular}{|p{1.8cm}|p{3.9cm}|p{2.9cm}|p{5.2cm}|}
\hline
\textbf{Prototipo} & \textbf{Directorio y tienda} & \textbf{Patrón y uso del artefacto} & \textbf{Papel metodológico} \\
\hline
\textbf{Prototipo A} & \texttt{checkout-con-sdk-\allowbreak solo-back/} («Austral Mate», puerto \texttt{4200}) & Patrón P1 parcial: \texttt{KitPagos} en el servidor; \texttt{fetch} nativo en el navegador. & Demuestra la adopción del SDK únicamente en el backend cuando el comercio conserva su propia captura en el frontend (fuera del experimento comparativo, puntos 123 y 124). \\
\hline
\textbf{Prototipo B} & \texttt{checkout-con-sdk-\allowbreak front-back/} («La Bodega de Chapinero», puerto \texttt{4101}) & Patrón P1 completo: \texttt{KitPagos} en el servidor y \texttt{KitPagosBrowser} en el navegador. & \textbf{Grupo experimental de tratamiento}: integra las 4 pasarelas mediante el SDK en backend y frontend (puntos 116 a 119 y 124). \\
\hline
\textbf{Prototipo C} & \texttt{checkout-frontend-\allowbreak rest/} («Tostadero Ceiba», puerto \texttt{5174}) & Patrón P2: cliente web estático que consume \texttt{/v1/api/*} y \texttt{KitPagosBrowser}. & Demuestra la viabilidad del contrato REST unificado desde un frontend sin servidor de pagos propio (fuera del experimento comparativo, puntos 91 a 106 y 124). \\
\hline
\textbf{Prototipo D} & \texttt{checkout-sin-\allowbreak solucion/} («La Candelaria Cómics», puerto \texttt{4102}) & Línea base manual: integración directa con \texttt{fetch} siguiendo la documentación oficial de cada pasarela. & \textbf{Grupo experimental de control}: integra las 4 pasarelas a mano sin usar el SDK ni la API de Simulación (puntos 110, 111, 117, 119 y 124). \\
\hline
\end{tabular}
\end{table}

El \textbf{Prototipo A} (\texttt{prototypes/checkout-con-sdk-solo-back/}, tienda «Austral Mate») representa el escenario de un comercio que adopta el paquete \texttt{kit-pagos-colombia} (\texttt{0.7.0-rc.2}) en su servidor Express (\texttt{src/PaymentService.ts}) para unificar la creación de pagos, la consulta de estados, el listado de bancos PSE y la verificación de webhooks a través de \texttt{KitPagos}, pero decide no incorporar \texttt{kit-pagos-colombia/browser} en su cliente web (\texttt{public/app.js}). En consecuencia, mientras su backend delega las cuatro pasarelas en una sola instancia del SDK por operación, su navegador construye manualmente las peticiones \texttt{fetch} de tokenización hacia Wompi (\texttt{POST /tokens/cards}), Mercado Pago (\texttt{POST /v1/card\_tokens}) y Kushki (\texttt{POST /card/v1/tokens}). Este prototipo documenta con código ejecutable y diez pruebas automatizadas el primer escalón de adopción del patrón P1 sin formar parte de la comparación experimental controlada (puntos 123 y 124).

Los \textbf{Prototipos B y D} conforman el par experimental diseñado para contrastar el desarrollo asistido por el artefacto frente a la integración tradicional en la Fase 5 de DSR (Peffers et al., 2007; puntos 110 a 124). Para aislar como única variable independiente el mecanismo de integración con las pasarelas, ambos prototipos comparten un esqueleto idéntico de servidor Express (\texttt{src/app.ts}, \texttt{src/server.ts}, \texttt{src/catalog.ts}, \texttt{src/OrderStore.ts}, \texttt{src/env.ts}, \texttt{src/trace.ts} y \texttt{src/types.ts}) y la misma estructura de interfaz de usuario en \texttt{public/}, cuya paridad fuera de la costura de pagos es verificada en cada ejecución de la CI por \texttt{prototypes/\allowbreak check-skeleton-parity.sh} (punto 116). Siguiendo el principio de aislamiento mediante puertos de la arquitectura hexagonal (Cockburn, 2024), la frontera entre el esqueleto común y la lógica de pagos está definida por el contrato \texttt{PaymentsPort} (\texttt{src/payments/PaymentsPort.ts}), compuesto por seis funciones: \texttt{listGateways}, \texttt{listPseBanks}, \texttt{createCardCheckout}, \texttt{createPseCheckout}, \texttt{refreshOrderStatus} y \texttt{handleWebhook}.

En el \textbf{Prototipo B} (\texttt{prototypes/checkout-con-sdk-front-back/}, tienda «La Bodega de Chapinero», grupo de tratamiento), la implementación de \texttt{PaymentsPort} reside en un único módulo de servidor (\texttt{src/payments/sdkAdapter.ts}) que instancia \texttt{KitPagos} con las credenciales de las cuatro pasarelas y resuelve las operaciones invocando \texttt{getPseBanks()}, \texttt{createPayment()}, \texttt{getPaymentStatus()} y \texttt{verifyWebhook()}, mientras que el navegador carga \texttt{/vendor/kit-pagos-browser.js} y tokeniza las tarjetas de Wompi, Mercado Pago y Kushki con una sola llamada a \texttt{KitPagosBrowser.tokenizeCard()} (puntos 116 y 118). Por el contrario, en el \textbf{Prototipo D} (\texttt{prototypes/checkout-sin-solucion/}, tienda «La Candelaria Cómics», grupo de control), esa misma costura (\texttt{src/payments/manualAdapter.ts}) exige cuatro módulos específicos por proveedor en el servidor (\texttt{wompiClient.ts}, \texttt{mercadopagoClient.ts}, \texttt{kushkiClient.ts} y \texttt{rapydClient.ts}) que implementan desde cero las firmas SHA-256 y HMAC-SHA256, el sondeo de \texttt{async\_payment\_url}, la búsqueda de BIN, el desglose de IVA, las cascadas de consulta de estado, la creación previa de clientes en Rapyd y los verificadores de webhook de cada pasarela siguiendo la documentación oficial de cada proveedor (Kushki, 2024; Mercado Pago, 2024; Rapyd Financial Technology, 2024; Wompi Colombia, 2024), además de tres bloques manuales de tokenización con \texttt{fetch} en el navegador (puntos 110, 111 y 117).

Tanto el Prototipo B como el Prototipo D se verificaron de extremo a extremo contra los cuatro sandboxes reales y contra la API de Simulación mediante el interruptor \texttt{PAYMENT\_TARGET} (\texttt{sandbox} o \texttt{simulator}). Durante los ensayos de retorno bancario y de notificaciones en vivo se resolvieron dos restricciones operativas de los proveedores que quedaron incorporadas en el esqueleto compartido: primero, dado que el validador de órdenes PSE de Mercado Pago rechaza \texttt{http://localhost} como dominio de retorno (\texttt{return\_url}), \texttt{src/env.ts} construye por omisión el origen local con \texttt{http://localtest.me::port} —dominio público que resuelve en DNS a \texttt{127.0.0.1} pero supera la validación de nombres de host del proveedor— (punto 120); segundo, para recibir webhooks reales firmados por Wompi, Mercado Pago y Rapyd en las tiendas locales, se integró el túnel HTTPS de dominio fijo en ngrok (\texttt{PUBLIC\_BASE\_URL}) y se documentó la extracción del secreto de aplicación del vendedor de prueba en Mercado Pago (puntos 120 y 121). Las mediciones cuantitativas de tamaño, complejidad ciclomática, acoplamiento y defectos obtenidas al comparar los Prototipos B y D pertenecen a la Fase 5 y se presentan en el capítulo VII.

Por último, el \textbf{Prototipo C} (\texttt{prototypes/checkout-frontend-rest/}, tienda «Tostadero Ceiba») demuestra el patrón de integración P2 (puntos 91 a 106 y 124). A diferencia de los otros tres prototipos, carece de backend de negocio propio: es una aplicación estática compilada con Vite y servida por un servidor de archivos estáticos de 58 líneas (\texttt{server.mjs}) que se comunica desde el navegador con la fachada \texttt{/v1/api/*} de la API de Simulación (ya sea en \texttt{http://localhost:3000} o en \texttt{https://kit-pagos-colombia.onrender.com}). El cliente consulta al iniciar \texttt{GET /v1/api/gateways} y adapta dinámicamente el formulario según el campo \texttt{cardCapture}: para Wompi, Mercado Pago y Kushki (\texttt{BROWSER\_TOKEN}), carga \texttt{kit-pagos-colombia/browser} desde jsDelivr mediante un \texttt{importmap} con verificación de integridad de subrecursos SHA-384 (\textit{Subresource Integrity}, SRI, punto 91), tokeniza la tarjeta directamente contra el proveedor y envía el token resultante a \texttt{POST /v1/api/payments}; para Rapyd (\texttt{HOSTED\_PAGE}), oculta los campos de tarjeta en la tienda, solicita la creación del cobro a \texttt{POST /v1/api/payments} y redirige el navegador a la página alojada devuelta en \texttt{redirectUrl} (punto 106). Como se advierte en el punto 114 de la bitácora de arquitectura, este diseño ilustra la ergonomía del contrato REST unificado, pero en un despliegue productivo de comercio electrónico la invocación a \texttt{POST /v1/api/payments} debe ejecutarse desde el servidor del comercio para que el monto de la orden no pueda ser manipulado desde las herramientas de desarrollo del navegador.

\fussy
